
\Needspace{15\baselineskip}
\subsubsection{Internal evaluations of explicit firms and posterior contrast}

The following table presents in legible form the seventeen HMT--MD coordinates already calculated---eleven elementary and six composite---and compares them with the 2026 edition of the Particle Data Group data set, whose editorial cutoff date is 15 January 2026 \cite{PDG2026}. They do not constitute an ontology distinct from the preceding operatorial layer. The symbol \(\dagger\) identifies an evaluation obtained by applying the construction map to an explicit signature. In the eleven elementary rows, the signature follows from the nominal route determined without masses; in the six composite rows, it follows from the declared binding ledger. Each figure depends on that signature, the bidirectional reader, and the common dimensional chart; the external data set enters only in the comparison and selects neither flavor, generation, color, nor route.


For charged leptons and massive bosons, the relative difference between
the two central coordinates is shown and, when the dataset provides a two-sided
uncertainty, its ratio to that uncertainty. For quarks, the
arithmetic difference between central values is reported, but the term
\emph{metrological residual} is reserved: it requires the internal mass and the external reference
to use the same definition, renormalization scheme and scale.
The light masses are read in the \(\overline{\mathrm{MS}}\) convention at
\(2\,\mathrm{GeV}\); those of \(c\) and \(b\), in
\(\mu=\overline m_q\); the top mass comes from kinematic reconstruction.

\begingroup
\setlength{\tabcolsep}{4.5pt}
\renewcommand{\arraystretch}{1.28}
\begin{longtable}{@{}>{\raggedright\arraybackslash}p{.12\textwidth}>{\raggedright\arraybackslash}p{.27\textwidth}>{\raggedright\arraybackslash}p{.27\textwidth}>{\raggedright\arraybackslash}p{.28\textwidth}@{}}
\caption{Internal evaluations of explicit firms and subsequent metrological contrast.}\label{tab:masas-evaluaciones-firmas-explicitas}\\
\toprule
\textbf{Status} & \textbf{Signature coordinate} & \textbf{External physical reference} & \textbf{Scope} \\
\midrule
\endfirsthead
\toprule
\textbf{State} & \textbf{Signature coordinate} & \textbf{External physical reference} & \textbf{Scope} \\
\midrule
\endhead
\midrule
\multicolumn{4}{r}{\footnotesize Continues on the following page.}\\
\endfoot
\bottomrule
\endlastfoot
% HMT-MASS-EXPLICIT-SIGNATURE:SM-LC-2
$\mu^-/\mu^+$ & $105.658360\allowbreak{}294435\allowbreak{}829303\allowbreak{}4\;\mathrm{MeV}/c^2$\textsuperscript{\(\dagger\)}; E3 tower, signature (42,-3,8,0,2) & $105.6583755\pm0.0000023\;\mathrm{MeV}/c^2$; leptonic rest mass; PDG average and declared metrological conversion & Relative difference: $-0.143912\allowbreak{}530\;\mathrm{ppm}$\textsuperscript{\(\dagger\)}; $-6.611114\,\sigma$. The HMT figure is the exact evaluation of the declared signature. \\
% HMT-MASS-EXPLICIT-SIGNATURE:SM-LC-3
$\tau^-/\tau^+$ & $1776.860917\allowbreak{}631432\allowbreak{}295595\;\mathrm{MeV}/c^2$\textsuperscript{\(\dagger\)}; E3 tower, signature (64,0,25,-1,6) & $1776.93\pm0.09\;\mathrm{MeV}/c^2$; direct PDG average & Relative difference: $-38.877371\allowbreak{}966\;\mathrm{ppm}$\textsuperscript{\(\dagger\)}; $-0.813547\,\sigma$. The HMT figure is the exact evaluation of the declared signature. \\
%% HMT-MASS-EXPLICIT-SIGNATURE:SM-Q-U1
$u/\bar u$ & $2.165924\allowbreak{}536173\allowbreak{}907\;\mathrm{MeV}/c^2$\textsuperscript{\(\dagger\)}; angular chart, nominal route signature (11,7) & $2.16\pm0.07\;\mathrm{MeV}/c^2$; mass $\overline{\mathrm{MS}}$ to $\mu=2\,\mathrm{GeV}$; 90\% C.L. & Difference between central values: $2742.840821\allowbreak{}253\;\mathrm{ppm}$\textsuperscript{\(\dagger\)}. It is not a metrological residual until the same scheme and the same scale have been fixed. \\
% HMT-MASS-EXPLICIT-SIGNATURE:SM-Q-D1 record
$d/\bar d$ & $4.679550\allowbreak{}162948\allowbreak{}874\;\mathrm{MeV}/c^2$\textsuperscript{\(\dagger\)}; chart angular, firma nominal of route (17,8) & $4.70\pm0.07\;\mathrm{MeV}/c^2$; mass $\overline{\mathrm{MS}}$ to $\mu=2\,\mathrm{GeV}$; 90\% C.L. & Difference between central values: $-4351.029159\allowbreak{}814\;\mathrm{ppm}$\textsuperscript{\(\dagger\)}. It is not a metrological residual until the same scheme and the same scale have been fixed. \\
%% HMT-MASS-EXPLICIT-SIGNATURE:SM-Q-U2
$c/\bar c$ & $1.270500\allowbreak{}838641\allowbreak{}911\;\mathrm{GeV}/c^2$\textsuperscript{\(\dagger\)}; angular chart, nominal route signature (61,8) & $1.2729\pm0.0045\;\mathrm{GeV}/c^2$; mass $\overline{\mathrm{MS}}$ evaluated at $\mu=\overline{m}_q$; 90\% C.L. & Difference between central values: $-1884.799558\allowbreak{}558\;\mathrm{ppm}$\textsuperscript{\(\dagger\)}. It is not a metrological residual until the same scheme and the same scale have been fixed. \\
%% HMT-MASS-EXPLICIT-SIGNATURE:SM-Q-D2
$s/\bar s$ & $92.940093\allowbreak{}907845\allowbreak{}64\;\mathrm{MeV}/c^2$\textsuperscript{\(\dagger\)}; angular chart, nominal route signature (41,-3) & $92.9\pm0.7\;\mathrm{MeV}/c^2$; mass $\overline{\mathrm{MS}}$ to $\mu=2\,\mathrm{GeV}$; 90\% C.L. & Difference between central values: $431.581354\allowbreak{}635\;\mathrm{ppm}$\textsuperscript{\(\dagger\)}. It is not a metrological residual until the same scheme and the same scale have been fixed. \\
% HMT-MASS-EXPLICIT-SIGNATURE:SM-Q-U3 record
$t/\bar t$ & $172.597479\allowbreak{}544019\allowbreak{}1\;\mathrm{GeV}/c^2$\textsuperscript{\(\dagger\)}; chart angular, firma nominal of route (100,-1) & $172.60\pm0.27\;\mathrm{GeV}/c^2$; mass reconstructed directly from event kinematics & Difference between central values: $-14.602873\allowbreak{}585\;\mathrm{ppm}$\textsuperscript{\(\dagger\)}. It is not a metrological residual until the same scheme and the same scale have been fixed. \\
%% HMT-MASS-EXPLICIT-SIGNATURE:SM-Q-D3
$b/\bar b$ & $4.190119\allowbreak{}763299\allowbreak{}790\;\mathrm{GeV}/c^2$\textsuperscript{\(\dagger\)}; angular chart, nominal route signature (71,-5) & $4.186\pm0.006\;\mathrm{GeV}/c^2$; mass $\overline{\mathrm{MS}}$ evaluated at $\mu=\overline{m}_q$; 90\% C.L. & Difference between central values: $984.176612\allowbreak{}467\;\mathrm{ppm}$\textsuperscript{\(\dagger\)}. It is not a metrological residual until the same scheme and the same scale have been fixed. \\
%% HMT-MASS-EXPLICIT-SIGNATURE:SM-GA-W
$W^\pm$ & $80.378964\allowbreak{}909704\allowbreak{}547024\allowbreak{}86\;\mathrm{GeV}/c^2$\textsuperscript{\(\dagger\)}; E3 tower, signature (94,-1,0,-2,4) & $80.3625\pm0.0077\;\mathrm{GeV}/c^2$; mass average of the sealed PDG snapshot & Relative difference: $204.882995\allowbreak{}234\;\mathrm{ppm}$\textsuperscript{\(\dagger\)}; $2.138299\,\sigma$. The HMT figure is the exact evaluation of the declared signature. \\
% HMT-MASS-EXPLICIT-SIGNATURE:SM-GA-Z record
$Z^0$ & $91.187533\allowbreak{}444313\allowbreak{}407844\allowbreak{}85\;\mathrm{GeV}/c^2$\textsuperscript{\(\dagger\)}; tower E3, signature (95,-1,-11,0,-4) & $91.1879\pm0.0020\;\mathrm{GeV}/c^2$; mass average of the sealed PDG snapshot & Relative difference: $-4.019784\allowbreak{}276\;\mathrm{ppm}$\textsuperscript{\(\dagger\)}; $-0.186362\,\sigma$. The HMT figure is the exact evaluation of the declared signature. \\
%% HMT-MASS-EXPLICIT-SIGNATURE:SM-H-1
Higgs $H$ & $125.292466\allowbreak{}042536\allowbreak{}133\;\mathrm{GeV}/c^2$\textsuperscript{\(\dagger\)}; angular chart, nominal route signature (97,9) & $125.13\pm0.11\;\mathrm{GeV}/c^2$; mass average of the sealed PDG snapshot & Relative difference: $1298.378027\allowbreak{}140\;\mathrm{ppm}$\textsuperscript{\(\dagger\)}; $1.454206\,\sigma$. The HMT figure is the exact evaluation of the declared signature. \\
% HMT-MASS-EXPLICIT-SIGNATURE:E3B-002
$\pi^0$\par signature $(44,\allowbreak -4,\allowbreak -25,\allowbreak 1,\allowbreak -2)$ & $134.976724\allowbreak{}327872\allowbreak{}125739\allowbreak{}384878\ldots\;\mathrm{MeV}/c^2$\textsuperscript{\(\dagger\)} & $134.976827\allowbreak{}767684\allowbreak{}71\pm 0.000485\allowbreak{}679282\allowbreak{}284233\allowbreak{}51\;\mathrm{MeV}/c^2$\par\footnotesize PDG 2026 catalog. & Relative difference $-0.766352\allowbreak{}375\,\mathrm{ppm}$; $-0.212979\,\sigma$. For the declared composite signature, the evaluation is deduced exactly from the beta E3 character. The prospective selector of that signature is an earlier map independent of the numerical comparison. \\
% HMT-MASS-EXPLICIT-SIGNATURE:E3B-003
$\pi^\pm$\par signature $(44,\allowbreak 1,\allowbreak -2,\allowbreak 2,\allowbreak -1)$ & $139.570485\allowbreak{}171572\allowbreak{}191374\allowbreak{}734364\ldots\;\mathrm{MeV}/c^2$\textsuperscript{\(\dagger\)} & $139.570390\allowbreak{}983681\allowbreak{}31\pm 0.000182\allowbreak{}007160\allowbreak{}408260\allowbreak{}09\;\mathrm{MeV}/c^2$\par\footnotesize PDG 2026 catalog. & Relative difference $0.674841\allowbreak{}491\,\mathrm{ppm}$; $0.517495\,\sigma$. For the declared composite signature, the evaluation is deduced exactly from the beta E3 character. The prospective selector of that signature is an earlier map independent of the numerical comparison. \\
% HMT-MASS-EXPLICIT-SIGNATURE:E3B-004
$K^\pm$\par signature $(54,\allowbreak -1,\allowbreak 17,\allowbreak -2,\allowbreak 2)$ & $493.677085\allowbreak{}649530\allowbreak{}612475\allowbreak{}723054\ldots\;\mathrm{MeV}/c^2$\textsuperscript{\(\dagger\)} & $493.676599\allowbreak{}458040\allowbreak{}58\pm 0.015405\allowbreak{}665226\allowbreak{}071509\;\mathrm{MeV}/c^2$\par\footnotesize PDG 2026 catalog. & Relative difference $0.984838\allowbreak{}030\,\mathrm{ppm}$; $0.031559\,\sigma$. For the declared composite signature, the evaluation is deduced exactly from the beta E3 character. The prospective selector of that signature is an earlier map independent of the numerical comparison. \\
% HMT-MASS-EXPLICIT-SIGNATURE:E3B-005
$K^0$\par signature $(54,\allowbreak 0,\allowbreak 32,\allowbreak 3,\allowbreak -2)$ & $497.611186\allowbreak{}497750\allowbreak{}457358\allowbreak{}558672\ldots\;\mathrm{MeV}/c^2$\textsuperscript{\(\dagger\)} & $497.610767\allowbreak{}531257\allowbreak{}52\pm 0.012990\allowbreak{}199756\allowbreak{}4242\;\mathrm{MeV}/c^2$\par\footnotesize PDG 2026 catalog. & Relative difference $0.841956\allowbreak{}244\,\mathrm{ppm}$; $0.032252\,\sigma$. For the declared composite signature, the evaluation is deduced exactly from the beta E3 character. The prospective selector of that signature is an earlier map independent of the numerical comparison. \\
% HMT-MASS-EXPLICIT-SIGNATURE:E3B-006
proton $p$\par signature $(59,\allowbreak 0,\allowbreak 9,\allowbreak 1,\allowbreak 0)$ & $938.271751\allowbreak{}546984\allowbreak{}898298\allowbreak{}936787\ldots\;\mathrm{MeV}/c^2$\textsuperscript{\(\dagger\)} & $938.272089\allowbreak{}430000\allowbreak{}05\pm 0.000000\allowbreak{}290000\allowbreak{}000000\allowbreak{}00\;\mathrm{MeV}/c^2$\par\footnotesize PDG 2026 catalog. & Relative difference $-0.360111\allowbreak{}974\,\mathrm{ppm}$; $-1165.113845\,\sigma$. For the declared composite signature, the evaluation is deduced exactly from the beta E3 character. The prospective selector of that signature is an earlier map independent of the numerical comparison. \\
% HMT-MASS-EXPLICIT-SIGNATURE:E3B-007
neutron $n$\par signature $(59,\allowbreak 1,\allowbreak -34,\allowbreak 0,\allowbreak 1)$ & $939.565810\allowbreak{}472507\allowbreak{}023312\allowbreak{}664431\ldots\;\mathrm{MeV}/c^2$\textsuperscript{\(\dagger\)} & $939.565421\allowbreak{}939999\allowbreak{}96\pm 0.000000\allowbreak{}480000\allowbreak{}000000\allowbreak{}00\;\mathrm{MeV}/c^2$\par\footnotesize PDG 2026 catalog. & Relative difference $0.413523\allowbreak{}633\,\mathrm{ppm}$; $809.442723\,\sigma$. For the declared composite signature, the evaluation is deduced exactly from the beta E3 character. The prospective selector of that signature is an earlier map independent of the numerical comparison. \\
\end{longtable}
\endgroup


\Needspace{12\baselineskip}
\paragraph{Native Invariants of Topological and Virtual Sectors.}
The Fibonacci, Ising, and \(\mathbb Z/6\) sectors and the virtual sections are compared in their proper codomains. Fibonacci publishes the based ring and its Frobenius--Perron dimension; the data \(F/R\) remain in the subsequent realization. Ising preserves its fusion ring and separates the Majorana realization; \(\mathbb Z/6\) publishes its braiding characters; virtuality publishes persistence horizons. This table maintains those invariants within the same realization map without artificially translating them into a unit of mass.

\begingroup
\setlength{\tabcolsep}{4.5pt}
\renewcommand{\arraystretch}{1.28}
\begin{longtable}{@{}>{\raggedright\arraybackslash}p{.15\textwidth}>{\raggedright\arraybackslash}p{.31\textwidth}>{\raggedright\arraybackslash}p{.25\textwidth}>{\raggedright\arraybackslash}p{.23\textwidth}@{}}
\caption{Topological and virtual sectors: native invariants of fusion, braiding, and persistence.}\label{tab:masas-invariantes-topologicos-virtuales-rev3}\\
\toprule
\textbf{Sector designation} & \textbf{Object and relations} & \textbf{Native invariant} & \textbf{Physical scope} \\
\midrule
\endfirsthead
\toprule
\textbf{Sector type} & \textbf{Object and relations} & \textbf{Native invariant} & \textbf{Physical scope} \\
\midrule
\endhead
\midrule
\multicolumn{4}{r}{\footnotesize Continues on the following page.}\\
\endfoot
\bottomrule
\endlastfoot
Fibonacci datum & $K=\mathbb Z1\oplus\mathbb Z\tau$, $\tau^2=1+\tau$. & $\operatorname{FPdim}(\tau)=\varphi$ in the based ring. & The matrices $F/R$, the pentagon, the hexagon, the Hamiltonian, and the gap belong to the subsequent realization. \\
Ising category & $K=\mathbb Z\{1,\psi,\sigma\}$, $\psi^2=1$, $\psi\sigma=\sigma$, $\sigma^2=1+\psi$. & $\operatorname{FPdim}(\sigma)=\sqrt{2}$; braiding representation. & The Majorana realization preserves its own physical morphism. \\
$\mathbb Z/6$ & $q_{\epsilon,t}=(-1)^\epsilon\omega_3^t=\zeta_6^k$. & Quadratic character and braid representation over the six sectors. & The relevant remark is topological, not a universal mass. \\
Secciones virtuales & Inverse system of prolongations with horizon $L(s_N)$. & Finite register, signature, carry, and obstruction class. & The native observable is the persistence of the prolongation. \\
\end{longtable}
\endgroup


Neutrinos remain in the codomain of eigenvalues and mixing described by their operator records; PMNS realizes transport between the interaction and mass bases, while the scale belongs to the spectrum of the neutrino operator. The six composite evaluations are obtained from explicit signatures through the general binding morphism into the scalar chart. The topological sectors and virtual sections retain their own fusion, braiding, and persistence observables.

\subsection{Causal construction and coverage of mass classes}

The APP--TRIT--TPK architecture, the mass law, the bidirectional reader, and the towers jointly determine the atlas, the enriched chronological ledgers, the signatures, the character, and the basal closure. Their confluence in \eqref{eq:ley-operatorial-general-masas-rectoras} provides a unique equation of mass sections. The complete domain is the atlas \(81/56/13/324\), and each realization specifies, according to its own codomain---that is, according to the codomain of each sector---the classifiers, cells, routes, orientations, dodecaphase words, signatures, and compatibility morphisms. The spectra \(K_D/K_\Omega\) refine the fibres of the atlas, while the chronological routes of \(108\) states associated with \(W\), \(Z\), and \(H\) directly construct their bosonic signatures and characters.
